% Worked Examples: Concept 1.2.3 - Stability
\documentclass[11pt,a4paper]{article}
\usepackage{worked-examples-template}

\title{\textbf{Worked Examples}\\[0.5em]
\large Concept 1.2.3: Stability\\[0.3em]
\normalsize \coursesubtitle{} -- \coursetitle}
\author{Assoc.\ Prof.\ William Robertson \and Dr Sean McGowan}
\date{\today}

\begin{document}

\maketitle
\tableofcontents
\newpage

\section{Concept 1.2.3: Stability}

\subsection{Example 1: Lyapunov Stability Analysis}

\paragraph{Problem:} Consider the nonlinear system:
\begin{align}
\dot{x}_1 &= -x_1 + x_1 x_2 \\
\dot{x}_2 &= -x_2
\end{align}

(a) Find the equilibrium points.

(b) Use a Lyapunov function candidate $V(x) = \frac{1}{2}(x_1^2 + x_2^2)$ to prove stability of the origin.

(c) Verify by linearization.

\paragraph{Solution:}

\textbf{Step 1: Equilibrium points}

Set $\dot{x}_1 = \dot{x}_2 = 0$:
\begin{align}
-x_1 + x_1 x_2 &= 0 \quad \Rightarrow \quad x_1(x_2 - 1) = 0 \\
-x_2 &= 0 \quad \Rightarrow \quad x_2 = 0
\end{align}

From the second equation, $x_2 = 0$. Substituting into the first: $x_1(-1) = 0$, so $x_1 = 0$.

The only equilibrium is the origin: $(0, 0)$.

\textbf{Step 2: Lyapunov stability analysis}

Consider $V(x) = \frac{1}{2}(x_1^2 + x_2^2)$. This is:
\begin{itemize}
\item Positive definite: $V(x) > 0$ for all $x \neq 0$ and $V(0) = 0$ ✓
\item Radially unbounded: $V(x) \to \infty$ as $\|x\| \to \infty$ ✓
\end{itemize}

Compute the time derivative along system trajectories:
\begin{align}
\dot{V} = \frac{\partial V}{\partial x_1}\dot{x}_1 + \frac{\partial V}{\partial x_2}\dot{x}_2
\end{align}

\begin{align}
\dot{V} &= x_1(-x_1 + x_1 x_2) + x_2(-x_2) \\
&= -x_1^2 + x_1^2 x_2 - x_2^2 \\
&= -x_1^2(1 - x_2) - x_2^2
\end{align}

\textbf{Analysis of $\dot{V}$:}
\begin{itemize}
\item If $|x_2| < 1$, then $1 - x_2 > 0$, so both terms are negative: $\dot{V} < 0$
\item If $x_2 > 1$, then $1 - x_2 < 0$, and the sign of $\dot{V}$ depends on relative magnitudes
\end{itemize}

For the region $|x_2| < 1$ (which includes a neighborhood of the origin), we have:
\begin{align}
\dot{V} \leq -\min(1 - x_2, 1) \cdot x_1^2 - x_2^2 < 0 \quad \text{for } x \neq 0
\end{align}

Since $V$ is positive definite and $\dot{V}$ is negative definite in a neighborhood of the origin, the origin is \textbf{asymptotically stable} by Lyapunov's theorem.

For global stability, we need $\dot{V} < 0$ everywhere. Rewrite:
\begin{align}
\dot{V} = -x_1^2 - x_2^2 + x_1^2 x_2 = -(x_1^2 + x_2^2) + x_1^2 x_2
\end{align}

For $|x_2| < 1$: $|x_1^2 x_2| < x_1^2 < x_1^2 + x_2^2$, so $\dot{V} < 0$ ✓

For $|x_2| \geq 1$: The analysis is more complex. Let's check a specific case: if $x_2 = 2$:
\begin{align}
\dot{V} = -x_1^2(1-2) - 4 = x_1^2 - 4
\end{align}
This can be positive if $|x_1| > 2$.

However, for the region near the origin (which is what matters for local stability), $\dot{V} < 0$ definitively.

\textbf{Conclusion:} The origin is \textbf{locally asymptotically stable}.

\textbf{Step 3: Verification by linearization}

The Jacobian at the origin is:
\begin{align}
J = \begin{bmatrix} \frac{\partial \dot{x}_1}{\partial x_1} & \frac{\partial \dot{x}_1}{\partial x_2} \\ \frac{\partial \dot{x}_2}{\partial x_1} & \frac{\partial \dot{x}_2}{\partial x_2} \end{bmatrix}\bigg|_{(0,0)} = \begin{bmatrix} -1 + x_2 & x_1 \\ 0 & -1 \end{bmatrix}\bigg|_{(0,0)} = \begin{bmatrix} -1 & 0 \\ 0 & -1 \end{bmatrix}
\end{align}

Eigenvalues: $\lambda_1 = \lambda_2 = -1$

Both eigenvalues have negative real parts $\Rightarrow$ the linearized system is stable $\Rightarrow$ the origin is \textbf{locally asymptotically stable} ✓

\textbf{Key Insights:}
\begin{itemize}
\item Lyapunov functions provide a powerful method to prove stability without solving the differential equations
\item If $V > 0$ and $\dot{V} < 0$ in a neighborhood, the equilibrium is locally asymptotically stable
\item If $V$ and $-\dot{V}$ are both positive definite and radially unbounded, the equilibrium is globally asymptotically stable
\item Finding Lyapunov functions is more art than science—common choices include quadratic forms, energy functions, or weighted sums
\item Linearization provides an alternative (often simpler) method for local stability analysis
\end{itemize}

\subsection{Example 2: Eigenvalue Stability Criterion}

\paragraph{Problem:} Determine the stability of the equilibrium point $(0, 0)$ for the following linear systems:

(a) $\begin{bmatrix} \dot{x}_1 \\ \dot{x}_2 \end{bmatrix} = \begin{bmatrix} -2 & 1 \\ 0 & -3 \end{bmatrix}\begin{bmatrix} x_1 \\ x_2 \end{bmatrix}$

(b) $\begin{bmatrix} \dot{x}_1 \\ \dot{x}_2 \end{bmatrix} = \begin{bmatrix} -1 & 4 \\ -1 & -1 \end{bmatrix}\begin{bmatrix} x_1 \\ x_2 \end{bmatrix}$

(c) $\begin{bmatrix} \dot{x}_1 \\ \dot{x}_2 \end{bmatrix} = \begin{bmatrix} 0 & 1 \\ -4 & 0 \end{bmatrix}\begin{bmatrix} x_1 \\ x_2 \end{bmatrix}$

\paragraph{Solution:}

\textbf{Stability criterion:} For a linear system $\dot{x} = Ax$, the origin is:
\begin{itemize}
\item \textbf{Asymptotically stable} if all eigenvalues of $A$ have negative real parts
\item \textbf{Unstable} if at least one eigenvalue has positive real part
\item \textbf{Marginally stable} if some eigenvalues are purely imaginary (and the rest have negative real parts) with simple multiplicity
\end{itemize}

\textbf{Part (a):}
\begin{align}
A = \begin{bmatrix} -2 & 1 \\ 0 & -3 \end{bmatrix}
\end{align}

Since $A$ is upper triangular, eigenvalues are on the diagonal: $\lambda_1 = -2$, $\lambda_2 = -3$

Both eigenvalues are negative $\Rightarrow$ \textbf{asymptotically stable}

The system is a stable node.

\textbf{Part (b):}
\begin{align}
A = \begin{bmatrix} -1 & 4 \\ -1 & -1 \end{bmatrix}
\end{align}

Characteristic equation:
\begin{align}
\det(sI - A) = \det\begin{bmatrix} s+1 & -4 \\ 1 & s+1 \end{bmatrix} = (s+1)^2 + 4 = s^2 + 2s + 5 = 0
\end{align}

\begin{align}
s = \frac{-2 \pm \sqrt{4-20}}{2} = \frac{-2 \pm 4j}{2} = -1 \pm 2j
\end{align}

Eigenvalues: $\lambda_1 = -1 + 2j$, $\lambda_2 = -1 - 2j$

Both have negative real parts ($-1$) $\Rightarrow$ \textbf{asymptotically stable}

The system is a stable spiral (oscillates while converging).

\textbf{Part (c):}
\begin{align}
A = \begin{bmatrix} 0 & 1 \\ -4 & 0 \end{bmatrix}
\end{align}

Characteristic equation:
\begin{align}
\det(sI - A) = \det\begin{bmatrix} s & -1 \\ 4 & s \end{bmatrix} = s^2 + 4 = 0
\end{align}

\begin{align}
s = \pm 2j
\end{align}

Eigenvalues: $\lambda_1 = 2j$, $\lambda_2 = -2j$ (purely imaginary)

Since eigenvalues are on the imaginary axis $\Rightarrow$ \textbf{marginally stable} (Lyapunov stable but not asymptotically stable)

The system exhibits sustained oscillations with constant amplitude (like an undamped oscillator).

\textbf{Key Insights:}
\begin{itemize}
\item Eigenvalues completely determine the stability of linear systems
\item The real part of eigenvalues determines growth/decay; the imaginary part determines oscillation frequency
\item Complex eigenvalues come in conjugate pairs for real matrices
\item Marginally stable systems are sensitive to perturbations and nonlinearities
\item In practice, systems should be designed with sufficient stability margin (eigenvalues well into the left half-plane)
\end{itemize}
\end{document}
